\begin{abstract}
  With deployment of standardized post-quantum cryptography
  well on the way, discussions continue around different approaches
  to hybrid deployments -- either combining post-quantum schemes
  with ``traditional'' elliptic-curve cryptography (or RSA) or 
  combining multiple different post-quantum schemes. 
  The typical question asked in these discussion is what the ``best''
  or ``most conservative'' approach is within a certain performance
  budget.  
  \vadim{because you mention Frodo later in the abstract, it would be good to say here that people also looked for schemes that eliminate potential points of weakness, such as Frodo, which aims to eliminate the algebraic structure.  I think the main motivation of our paper is that Frodo+Kaiburr is safer than Frodo, so we should state something about Frodo up front.}

  In this paper we add a data point to this discussion by investigating
  how to scale ML-KEM to (much) higher security levels of the underlying
  MLWE problem. We first show that increasing security 
  simply by increasing the module rank $k$ 
  quickly reaches a limit imposed by decryption failure probability.
  We then propose a generalization of the noise distribution, which enables scaling
  to much larger module rank and MLWE hardness, while preserving the essence
  of ML-KEM, most notably arithmetic in the ring $\Z_{3329}[X]/(X^{256}+1)$ 
  % \vadim{and simple noise generation}.

  We select three concrete parameter sets of this high-security variant
  of ML-KEM, which we call \textsc{Kaiburr}. We present benchmarking results
  for optimized implementations of \kaiburr on 64-bit Intel CPUs and on
  the Asymmetric Crypto Coprocessor (ACC) of the Pavona open-source silicon framework.
  Based on these benchmarks, we argue that, for example, 
  a hybrid of FrodoKEM-976 and \kaiburr-1792 not only offers better
  performance than FrodoKEM-1344, it also offers protection against
  a broader range of hypothetical future attacks against (M)LWE.
\end{abstract}
